\section*{अध्याय \ref{ch:b1:heat-budget} --- पृथ्वी का ऊष्मा-बजट}

\begin{solution}{exo:b1:heat-budget:1}
$q = k\,\mathrm dT/\mathrm dz = 2.5\times\qty{0.022}{K.m^{-1}} =
\qty{0.055}{W.m^{-2}} = \qty{55}{mW.m^{-2}}$। \qty{1}{km^2} $=
\qty{e6}{m^2}$ के माध्यम से: \qty{5.5e4}{W}, अर्थात् \qty{55}{kW}।
\end{solution}

\begin{solution}{exo:b1:heat-budget:2}
महाद्वीप $\num{2.073e14}\times0.0709 = \qty{1.470e13}{W} = \qty{14.70}{TW}$;
महासागर $\num{3.028e14}\times0.1054 = \qty{31.92}{TW}$; कुल \qty{46.61}{TW}, और
औसत \omterm{def:b1:heat-budget:heat-flow}{ऊष्मा-प्रवाह} $\num{4.661e13}/\num{5.101e14} = \qty{91.4}{mW.m^{-2}}$ (प्रकाशित
46.7 और 91.6 केवल दो माध्यों के पूर्णांकन से भिन्न हैं)। सौर आगत का अंश $46.6/173000 = \num{2.7e-4}$ है, या
\qty{0.027}{\percent}।
\end{solution}

\begin{solution}{exo:b1:heat-budget:3}
यूरेनियम: $0.99274\times\num{9.513e-5} + 0.00720\times\num{5.685e-4} =
\num{9.444e-5} + \num{4.09e-6} = \qty{9.85e-5}{W.kg^{-1}}$; दुर्लभ
\ce{^{235}U} उसका \qty{4}{\percent} देता है। पोटैशियम: \ce{^{40}K} का द्रव्यमान-अंश
$\num{1.17e-4}\times39.96/39.10 = \num{1.196e-4}$ है, अतः
$\num{1.196e-4}\times\num{2.847e-5} = \qty{3.40e-9}{W.kg^{-1}}$।
\end{solution}

\begin{solution}{exo:b1:heat-budget:4}
$\qty{1}{\micro cal.cm^{-2}.s^{-1}} = \num{4.184e-6}\,\unit{J.s^{-1}}/\qty{e-4}{m^2}
= \qty{41.84}{mW.m^{-2}}$: पूर्वी संयुक्त राज्य \qty{33.5}{mW.m^{-2}}, बेसिन और
रेंज \qty{58.6}{mW.m^{-2}}, सिएरा नेवादा \qty{16.7}{mW.m^{-2}}। $A_0D$ को \unit{mW.m^{-2}} में लेते हुए, जब $A_0$
\unit{\micro W.m^{-3}} में और $D$ \unit{km} में हो:
$33.5 + 3.0\times7.5 = \qty{56.0}{mW.m^{-2}}$, $58.6 + 3.0\times9.4 =
\qty{86.8}{mW.m^{-2}}$, $16.7 + 3.0\times10 = \qty{46.7}{mW.m^{-2}}$। वही ग्रेनाइट
बेसिन और रेंज में सिएरा नेवादा की तुलना में लगभग दोगुने \omterm{def:b1:heat-budget:heat-flow}{ऊष्मा-प्रवाह} देता है: अंतर नीचे के \omterm{def:b1:layered-earth:layers}{प्रावार} में है।
\end{solution}

\begin{solution}{exo:b1:heat-budget:5}
$T(10\,\text{km}) = 10 + \frac{0.060\times10^4}{2.5} - \frac{\num{2e-6}\times10^8}{2\times2.5}
= 10 + 240 - 40 = \qty{210}{\degreeCelsius}$। नीचे, $q_m = 60 - 2.0\times10 =
\qty{40}{mW.m^{-2}}$, \qty{16}{K.km^{-1}} की प्रवणता, और
$T(35\,\text{km}) = 210 + 16\times25 = \qty{610}{\degreeCelsius}$।
\end{solution}

\begin{solution}{exo:b1:heat-budget:6}
शेल: प्रवणता $12.0/0.300 = \qty{40}{K.km^{-1}}$, $q = 1.5\times0.040 =
\qty{60}{mW.m^{-2}}$। ग्रेनाइट: $10.0/0.500 = \qty{20}{K.km^{-1}}$,
$q = 3.0\times0.020 = \qty{60}{mW.m^{-2}}$। $R(300) = 300/1.5 =
\qty{200}{m^2.K.W^{-1}}$, और इसके लिए $R(800) = 200 + 500/3.0 =$ $\qty{366.7}{m^2.K.W^{-1}}$।
तीन बिंदु $(0, 12.0)$, $(200, 24.0)$, $(366.7, 34.0)$ सरेख हैं, जिनकी ढलान $22.0/366.7 = \qty{0.0600}{W.m^{-2}}$ है: \omterm{def:b1:heat-budget:heat-flow}{ऊष्मा-प्रवाह} दोनों चट्टानों में
समान है, \qty{60}{mW.m^{-2}}; केवल प्रवणता बदलती है, चालकता के व्युत्क्रम अनुपात में।
\end{solution}

\begin{solution}{exo:b1:heat-budget:7}
$\tau = L^2/\kappa$: \qty{1}{m}, \qty{e6}{s} या \num{12}~दिन; \qty{10}{m},
\qty{e8}{s} या \num{3.2}~वर्ष; \qty{100}{km}, \qty{e16}{s} या
\qty{317}{Myr}; \qty{2900}{km}, \qty{8.4e18}{s} या \qty{267}{Gyr}। \qty{10}{m} पर दैनिक चक्र बिल्कुल महसूस नहीं होता, और वार्षिक चक्र, जो \qty{10}{m} पार करने में लगने वाले तीन वर्षों से छोटा है, बहुत दुर्बलता से: केवल वर्षों या उससे अधिक तक चलने वाले परिवर्तन आते हैं। अंतिम समय $267/4.54 = 59$ गुना है पृथ्वी की आयु का।
\end{solution}

\begin{solution}{exo:b1:heat-budget:8}
(क) $1646 - 0.54\times50 = \qty{1619}{K}$, अर्थात् \qty{1346}{\degreeCelsius},
प्रकाशित \qty{1350}{\degreeCelsius} का पूर्णांकित मान। (ख) माध्य प्रवणता
$(0.54 + 0.36)/2 = \qty{0.45}{K.km^{-1}}$, अतः $1646 + 0.45\times360 =
\qty{1808}{K}$, \qty{1799}{K} से \qty{9}{K} (\qty{0.5}{\percent}) ऊपर: प्रवणता रैखिक रूप से नहीं घटती, पर अनुमान निकट है।
\end{solution}

\begin{solution}{exo:b1:heat-budget:9}
(क) $G^2 = T_0^2/(\pi\kappa t)$ से $t =$ $T_0^2/(\pi\kappa G^2)$।
(ख) $G = 0.02\times\tfrac59\,\unit{K}/\qty{0.3048}{m} =$ $\qty{0.03645}{K.m^{-1}} =$
$\qty{36.45}{K.km^{-1}}$; $t = 3900^2/(\pi\times\num{1.2e-6}\times0.03645^2) =
\num{1.521e7}/\num{5.009e-9} = \qty{3.04e15}{s}$, अर्थात् \qty{96}{Myr}, केल्विन का आँकड़ा।
(ग) अब $t = 1300^2/(\pi\times\num{e-6}\times0.036^2) =
\qty{4.15e14}{s} = \qty{13}{Myr}$: अधिक यथार्थवादी आरंभिक तापमान चालक पृथ्वी को और भी युवा बना देता है।
\end{solution}

\begin{solution}{exo:b1:heat-budget:10}
(क) क्षरण के बाद नई सतह पुरानी गहराई $e$ है, जहाँ ऊष्मा-उत्पादन $A_0' = A_0e^{-e/D}$ था और \omterm{def:b1:heat-budget:heat-flow}{ऊष्मा-प्रवाह} $q(e) = q_m +
A_0De^{-e/D} = q_m + A_0'D$ था; नई सतह से मापने पर उत्पादन अब भी $A_0'e^{-z'/D}$ के अनुसार घटता है। प्लूटॉन उसी रेखा पर नीचे खिसकता है।
(ख) परत अब $h - e$ मोटी है, उसी $A$ के साथ, अतः $q_0' = q_m + A(h - e)$: ढलान $h - e$ की एक रेखा पर, $q_m$ से होकर।
(ग) केवल घातांकी मॉडल: एकसमान-परत मॉडल में भिन्न गहराइयों से क्षरित प्लूटॉन भिन्न ढलानों की रेखाओं पर होते, और एक प्रांत के बिंदु बिखर जाते। कि एक ही रेखा देखी जाती है, यह गहराई के साथ लगभग घातांकी रूप से घटते उत्पादन के पक्ष में है।
\end{solution}

\begin{solution}{exo:b1:heat-budget:11}
(क) $\tau_P = RL/(3\kappa)$: \qty{50}{km} के लिए, $\num{6.4e6}\times\num{5e4}/\num{3e-6}
= \qty{1.07e17}{s} = \qty{3.4}{Gyr}$; \qty{100}{km} के लिए, \qty{6.8}{Gyr}।
(ख) $q = kT_i/L = 3\times1300/\num{5e4} = \qty{78}{mW.m^{-2}}$ आरंभ में,
और $78\,e^{-2/3.38} = \qty{43}{mW.m^{-2}}$ \qty{2}{Gyr} के बाद। (ग)
$3\times1300/\sqrt{\pi\times\num{e-6}\times\num{6.31e16}} = 3900/\num{4.45e5} =
\qty{8.8}{mW.m^{-2}}$: अच्छी तरह मिश्रित पृथ्वी दो अरब वर्ष बाद भी चालक पृथ्वी की तुलना में पाँच गुना अधिक ऊष्मा खोती है। उसका पूरा भीतरी भाग सतह को पोषित करता है, जहाँ चालक पृथ्वी में केवल एक बाहरी कवच करता है।
\end{solution}

\begin{solution}{exo:b1:heat-budget:12}
(क) $\sqrt{\pi\kappa\times\qty{1}{Myr}} = \sqrt{\pi\times\num{e-6}\times\num{3.156e13}}
= \qty{9957}{m}$, अतः $C = 3\times1336/9957 = \qty{0.403}{W.m^{-2}}$, $t$ \unit{Myr} में होने पर \qty{403}{mW.m^{-2}.Myr^{1/2}}: आँकड़ों पर फ़िट
490 का \qty{82}{\percent}। मॉडल को बड़े $k\,\Delta T$ की आवश्यकता है (लगभग \qty{4900}{W.m^{-1}}): \omterm{def:b1:layered-earth:layers}{प्रावार} की उच्च चालकता, या रुद्धोष्म की तरह गरम भीतरी भाग।
(ख) $490/\sqrt1 = 490$, $490/\sqrt{10} = 155$ और $490/\sqrt{100} =
\qty{49}{mW.m^{-2}}$। (ग) $z = 2.326\sqrt{\num{e-6}\times\num{3.156e15}} =
2.326\times\qty{56180}{m} = \qty{131}{km}$।
\end{solution}

\begin{solution}{pb:b1:heat-budget:1}
\textbf{1.} प्रवणता $= (T_{500} - T_{100})/\qty{0.4}{km}$: A \qty{15.25}{K.km^{-1}},
B \num{19.50}, C \num{19.50}, D \num{25.50}, E \qty{28.25}{K.km^{-1}}.

\textbf{2.} ड्रिलिंग चट्टान के तापमान को गड़बड़ाती है (ड्रिलिंग-द्रव दीवारों को ठंडा या गरम करता है), और गड़बड़ी को शांत होने में सप्ताह से महीनों लगते हैं। ऊपर की दसियों मीटर ऋतु-चक्र और सतह-तापमान के हाल के परिवर्तनों का अनुसरण करती हैं।

\textbf{3.} $q_0 = k\times$प्रवणता: A $2.9\times15.25 = \qty{44.225}{mW.m^{-2}}$,
B \qty{50.70}{mW.m^{-2}}, C \qty{60.45}{mW.m^{-2}}, D \qty{68.85}{mW.m^{-2}},
E \qty{84.75}{mW.m^{-2}}.

\textbf{4.} A से D महाद्वीपीय औसत से नीचे हैं, E उससे \qty{13.9}{mW.m^{-2}} ऊपर, या \qty{20}{\percent}।

\textbf{5.} \omterm{def:b1:heat-budget:heat-flow}{ऊष्मा-प्रवाह} चालकता गुना प्रवणता है: C बेहतर चालन करता है (3.1 बनाम 2.6), अतः वही प्रवणता \qty{19}{\percent} अधिक ऊष्मा ढोती है।

\textbf{6.} समतापीय रेखाएँ उच्चावच का केवल आंशिक अनुसरण करती हैं: किसी शिखर के नीचे वे सतह से कम ऊपर उठती हैं, अतः वे उसके नीचे अधिक दूर-दूर होती हैं, और मापी गई प्रवणता घट जाती है (घाटी के नीचे बढ़ जाती है)। प्लूटॉनों की तुलना से पहले उच्चावच के लिए सुधार आवश्यक है।

\textbf{7.} $\rho C_U H_U$: A $2650\times\num{2.0e-6}\times\num{9.85e-5} =
\qty{0.522}{\micro W.m^{-3}}$, B 0.914, C 1.305, D 1.827, E 2.610.

\textbf{8.} थोरियम: A 0.558, B 0.976, C 1.255, D 1.742, E 2.439.

\textbf{9.} पोटैशियम: A $2650\times0.030\times\num{3.40e-9} = 0.270$, B 0.315,
C 0.360, D 0.378, E 0.405.

\textbf{10.} $A_0$ = 1.35, 2.20, 2.92, 3.95 और \qty{5.46}{\micro W.m^{-3}}।
पोटैशियम सबसे कम योगदान देता है, A में कुल का \qty{20}{\percent} से E में \qty{7}{\percent} तक; यूरेनियम और थोरियम शेष को लगभग समान रूप से बाँटते हैं।

\textbf{11.} हाँ: पाँचों बिंदु लगातार चढ़ते हैं, एक सीधी रेखा से लगभग
\qty{1.5}{mW.m^{-2}} के भीतर।

\textbf{12.} $\bar A_0 = 15.88/5 = 3.176$, $\bar q_0 = 308.975/5 = 61.795$;
$\sum(A_0 - \bar A_0)^2 = 10.1681$; $\sum(A_0 - \bar A_0)(q_0 - \bar q_0) =
101.146$.

\textbf{13.} $D = 101.146/10.1681 = 9.95$ (\unit{mW.m^{-2}} प्रति
\unit{\micro W.m^{-3}}, अर्थात् km में): $D = \qty{9.95}{km}$, और $q_m = 61.795 -
9.95\times3.176 = \qty{30.20}{mW.m^{-2}}$।

\textbf{14.} ढलान, लगभग \qty{10}{km}, तीन प्रांतों की है; न्यूनित \omterm{def:b1:heat-budget:heat-flow}{ऊष्मा-प्रवाह}, \qty{30}{mW.m^{-2}}, पूर्वी संयुक्त राज्य के \qty{33}{mW.m^{-2}} के निकट है: एक शांत प्रांत जिसके नीचे ठंडा \omterm{def:b1:layered-earth:layers}{प्रावार} है, बेसिन और रेंज से भिन्न।

\textbf{15.} $q_m/k = \qty{12}{K.km^{-1}}$ और $A_0D^2/k = \num{5.46e-6}\times\num{e8}/2.5
= \qty{218.4}{K}$: $T(z) = 10 + 12z + 218.4\,(1 - e^{-z/10})$, $z$ km में, $T$
\unit{\degreeCelsius} में।

\textbf{16.} $q(z) = 30 + 54.6\,e^{-z/10}$: \qty{10}{km} पर \qty{50.1}{mW.m^{-2}},
\qty{20}{km} पर \num{37.4} और \qty{35}{km} पर \qty{31.7}{mW.m^{-2}}।

\textbf{17.} $30/84.6 = \qty{35}{\percent}$: E पर \omterm{def:b1:heat-budget:heat-flow}{ऊष्मा-प्रवाह} के लगभग दो-तिहाई भाग का उत्पादन उसकी रेडियोजनिक परत में होता है।

\textbf{18.} $T(10) = 10 + 120 + 218.4\times0.632 = \qty{268}{\degreeCelsius}$.

\textbf{19.} $T(20) = 10 + 240 + 218.4\times0.865 = \qty{439}{\degreeCelsius}$.

\textbf{20.} A के नीचे, $A_0D^2/k = 1.35\times100/2.5 = \qty{54.0}{K}$:
$T(35) = 10 + 420 + 54.0\times0.970 = \qty{482}{\degreeCelsius}$।

\textbf{21.} $A_0D^2/k$ से: A के नीचे \qty{54}{K} और E के नीचे \qty{218}{K}। गहरा \omterm{def:b1:heat-budget:heat-flow}{ऊष्मा-प्रवाह} एक ही है; E की \omterm{def:b1:layered-earth:layers}{भूपर्पटी} ऊष्मा-स्रोतों में चार गुना समृद्ध है, और वे जो ऊष्मा उत्पन्न करते हैं उसे उसी चट्टान से ऊपर चालित होना होगा, जिसके लिए पूरे रास्ते एक तीखी प्रवणता चाहिए।

\textbf{22.} $1300 = 10 + 218.4 + 12z$ से $z = \qty{89.3}{km}$ E के नीचे; $1300 = 10 + 54.0 + 12z$ से \qty{103.0}{km} A के नीचे।

\textbf{23.} दोनों महाद्वीपीय \omterm{def:b1:layered-earth:lithosphere}{स्थलमंडल} की परास में हैं, युवा क्षेत्रों के नीचे सौ किलोमीटर या उससे कम से लेकर क्रेटनों के नीचे लगभग \qty{250}{km} तक, उसकी पतली ओर के निकट: \qty{30}{mW.m^{-2}} के न्यूनित \omterm{def:b1:heat-budget:heat-flow}{ऊष्मा-प्रवाह} और इस चालकता के साथ, चालक \omterm{def:b1:heat-budget:geotherm}{जियोथर्म} लगभग \qty{100}{km} पर प्रावार-तापमान तक पहुँचता है; अधिक रेडियोधर्मी प्लूटॉन के नीचे, \qty{14}{km} पहले।

\textbf{24.} $T(35) = 10 + 420 + 218.4\times(1 - e^{-3.5}) = 430 + 211.8$: प्लूटॉन E के नीचे \omterm{def:b1:layered-earth:discontinuity}{मोहो} पर तापमान
\textbf{\qty[text-series-to-math=true]{642}{\degreeCelsius}}।
\end{solution}
